% =========================================================================
% $VII$ Floor-side design rule and its certification status (lane C).
% Numbers: .work3/c90b_out.txt (floor descent), community board \S2.4/\S2.5/\S43,
% .work3/c13_joint.py::Phi (independent implementation).
% =========================================================================
\section{The floor-side design rule, and what is certified}\label{sec:floor}

\S\ref{sec:infeasible} fixes the \emph{wall} $D_{\min}(V)=j_c+\Phi_0(V)$ and shows
it does not order the frontiers. Designing a sensor nevertheless means choosing
the subspace, and on the floor side that choice is a lower-dimensional problem:

\begin{equation}
  \begin{aligned}
  &\Phi_0^{\star}(r):=\min_{V\in\mathrm{Gr}(n,r)}\Phi_0(V),\\
  &\Phi_0(V)=\inf_{\tau>0}\operatorname{tr}\big(\Theta P_m(\tau P_V)\big),
  \end{aligned}
  \label{eq:flooropt}
\end{equation}
with $j_c=\operatorname{tr}(WP_c)$ and $\Theta=K^{\top}(R+B^{\top}P_cB)K$ as in
\eqref{eq:cost}. Its solution is a \emph{staircase}, and the rule we report is the
one that attains it on the published example.

\paragraph{Rule.} At rank $r$, take the $r$ leading eigenvectors of the
closed-loop task weight $\Theta$ --- i.e.\ the $r$ directions the cost penalises
most --- and put the sensor there. The resulting floors are
\begin{equation}
  \Phi_0^{\rm rule}=12.222448,\;0.791828,\;0.190335
  \quad(r=1,2,3),
  \label{eq:rule}
\end{equation}
i.e.\ $+38.8\%$, $+2.51\%$, $+0.60\%$ over $j_c$, against the best known floors
\begin{equation}
  \begin{aligned}
  &\Phi_0^{\star}=10.169224,\;0.735143,\;0.158701,\\
  &D_{\min}=41.652566,\;32.218485,\;31.642043,
  \end{aligned}
  \label{eq:bestfloor}
\end{equation}
found by direct descent on $\mathrm{Gr}(4,r)$. The rule is therefore
$1.20/1.08/1.20\times$ the best known floor at $r=1/2/3$: it is a usable rule and
it is not optimal, and the gap is largest exactly where the task is most
constrained.

\paragraph{Certification status (read this before quoting \eqref{eq:bestfloor}).}
The numbers in \eqref{eq:bestfloor} come from multi-start Nelder--Mead on
$\mathrm{Gr}(4,r)$, not from a proof. Three facts bound what they certify.
\emph{(a)} The landscape is two-valued at rank one, with the two values
$10.169224$ and $86.665888$ --- the good basin is real, but the descent is a
search.
\emph{(b)} Re-running the same descent with an independent implementation (this
lane, DARE-based, $24$ starts) reproduces the \emph{values} in
\eqref{eq:bestfloor} to all printed digits, while the number of starts reaching
the best basin differs between runs ($24/24$ in the first run reported on the
board versus $14/24$, $7/24$, $3/24$ at $r=1,2,3$ in our rerun). The value is
stable; the \emph{basin count is not a certificate}, and neither is any single
run's.
\emph{(c)} The three independent implementations of $\Phi_0$ available to us
(an information-form noiseless fixed point, a projection iteration, and a
DARE-based evaluation) agree to $1.3\times10^{-3}$ absolute over the tested
designs, which certifies the \emph{evaluation}, not the minimisation.
We therefore write ``best known, certified global \emph{on the floor side}'' and
never ``the optimal $F$''.

\paragraph{It is not a new design.} Choosing $F$ by \eqref{eq:flooropt} is
equivalent to the linear-class optimum of the SDP in \cite{tanaka2015joint} read
on the floor side: the floor-side rule selects the same subspace the SDP would
select, and it does so with a scalar functional of $V$ instead of a matrix
program. Our contribution here is the floor-side reading and the rule, not the
design.

\paragraph{Two axes, not one.} That a design is \emph{lossless on the rate side}
says nothing about its floor, and conversely. The four corners measured on the
published example are:
\begin{center}\footnotesize
\begin{tabular}{@{}lcc@{}}
\toprule
subspace $V$ & rate loss $\Delta$ (bit) & floor excess $\Phi_0$ \\
\midrule
full state                          & $\approx0$        & $\approx0$ \\
$\operatorname{span}(S^{\star}(80))$ & $-6.7\times10^{-9}$ & $+38\%$ \\
Schur tail, $k=2$                   & $+0.03$           & $+46.5\%$ \\
Schur tail, $k=3$                   & $+1.04$           & $+16.3\%$ \\
\bottomrule
\end{tabular}
\end{center}
so $\Delta$ and $\Phi_0$ are independent axes: a lossless representation can be
$38\%$ more expensive than the best floor, and a badly lossy one can be close to
it. The precise floor-side statement is Proposition~\ref{pr:nofactor}: $\Phi_0=0$
iff the stage cost sees only the task variable,
$\operatorname{supp}(Q)\subseteq\operatorname{range}(F^{\top})$; and
$\Phi_0<\infty$ iff the pair $(A,F)$ is detectable
(Proposition~\ref{pr:region}), which is the reference paper's own Theorem~2.1 and
is quoted, not claimed. The dual reading is worth stating: detectability is what
keeps the rate side finite, and the same condition keeps the cost side finite, so
neither floor can be bought with the other.

\paragraph{Sensitivity.} Both of the above are deterministic functionals of the
plant, so finite-sample effects enter only through estimation error. Under a
relative perturbation $\varepsilon\le1\%$ the rule keeps the staircase ordering and
magnitude; beyond about $\varepsilon\gtrsim3\%$ a \emph{frozen} rule degrades
sharply (median $+28\%$) and must be re-derived. The absolute floor is dominated
by $j_c$ (sensitivity $\sim\varepsilon$), whereas staircase statements inherit the
$\sim10\varepsilon$ sensitivity of the floor functional; we therefore attach the
error level to every staircase number.
